\documentclass[11pt,a4paper]{article}
\usepackage[margin=19mm,headheight=15pt]{geometry}
\usepackage{fontspec}
\setmainfont{Noto Serif}
\usepackage{amsmath,amssymb,mathtools}
\usepackage{xcolor}
\usepackage{fancyhdr}
\usepackage{enumitem}
\usepackage{microtype}
\setlength{\parindent}{0pt}
\setlength{\parskip}{6pt}
\definecolor{navy}{HTML}{183153}
\definecolor{lightgray}{HTML}{F3F5F7}
\pagestyle{fancy}\fancyhf{}
\lhead{\small ME 5309 -- Mid-Semester Examination}
\rhead{\small ME 5309}
\cfoot{\thepage}
\begin{document}
\begin{center}
{\Large\bfseries ME 5309 -- Mid-Semester Examination}\\[3pt]
{\large Answer Booklet}\\[8pt]
\end{center}
\textbf{Instructions:} Attempt any three questions. Each question carries 10 marks. Show working where appropriate.
\section*{Question 1 \hfill [10 marks]}
\subsection*{Q1(a) \hfill [4 marks]}

Pixel pitch = 0.005 mm, sensor width = 1280x0.005 = 6.4 mm. how to calculate the belt field of view.

\subsection*{Q1(b) \hfill [3 marks]}

\[\begin{aligned}Sensitivity = 10 mV/N\end{aligned}\]
The ADC has 12 bits, so there are 4096 levels.

\subsection*{Q1(c) \hfill [3 marks]}

\[\begin{aligned}F = e^(AT) = e^(-0.2) = 0.818731\end{aligned}\]

\section*{Question 3 \hfill [10 marks]}
\subsection*{Q3(a) \hfill [2 marks]}

The estimate is noisy and can fluctuate from sample to sample.

\subsection*{Q3(b) \hfill [2 marks]}

Tustin uses s = (2/T)(z-1)/(z+1).

\subsection*{Q3(c) \hfill [4 marks]}

A-LC=[[-ℓ₁,1], [-ℓ₂,-2]].
The desired polynomial is (s+10)² = s²+20s+100.

\subsection*{Q3(d) \hfill [2 marks]}

An observer uses the plant model and known control input to estimate unmeasured states such as velocity, and its poles can be designed.
A simple low-pass filter cannot exploit the state-space model in this way.
A filter may be safer when the model is poorly known or unreliable, because it relies less on model accuracy.\par\smallskip The observer uses the plant model and the measured output to reconstruct the unmeasured state. Its convergence depends on the model and the observer poles. A low-pass filter is simpler, but generally introduces phase lag and does not explicitly use the plant dynamics.

\section*{Question 4 \hfill [10 marks]}
\subsection*{Q4(a) \hfill [2 marks]}

The characteristic polynomial is s²+4s+8, so I report poles -4 and -8.
Unity DC gain requires Nff=8.

\subsection*{Q4(b) \hfill [2 marks]}

The actual feedforward contribution is 0.75x8=6.

\subsection*{Q4(c) \hfill [4 marks]}

\[\begin{aligned}k_1=18, k_2=6, K_I=32\end{aligned}\]

\subsection*{Q4(d) \hfill [2 marks]}

\vspace{2mm}

\end{document}